\section{Conclusion}
\label{sec:conclusion}

We opened with a structural paradox: the best-performing method for spurious correlation robustness achieves WGA = 0.91 without changing the backbone at all. Our diagnostic study resolves this paradox not as an anomaly but as a structural feature of the robustification landscape. WGA improvement comes from two mechanistically distinct pathways --- head recalibration on unchanged backbone features (DFR), and backbone-level spurious encoding reduction (GroupDRO) --- and both are effective in the Waterbirds WILDS setting.

For GroupDRO, we verify a three-step mechanistic chain: minority group upweighting (5.01\% of training data, exponentiated gradient ascent) $\to$ group-balanced gradient propagating through all backbone layers (backbone/head L2 ratio = 6.47--7.03) $\to$ significantly reduced background linear decodability in layer4 features (ERM 0.9838 $\to$ GroupDRO 0.9530, $p = 0.0039$, Cohen's $d = 6.48$, CONFIRMED). For DFR, we provide the first quantitative confirmation that its backbone is numerically identical to ERM at matching seeds (cosine similarity = 1.000000 $\pm$ 1e-14).

The backbone-vs-head typology matters because it changes how we think about method selection, transfer learning, and future method design. The probe protocol we establish --- frozen layer4 features, sklearn L-BFGS $C=10^9$, full test set, paired t-test across matched seeds --- is a low-cost, reproducible backbone spurious encoding diagnostic.

What remains open: the suppression-vs-dilution mechanistic ambiguity requires directional subspace analysis; the WGA correlation needs per-seed WGA measurements for full confirmation; layer-wise spurious encoding profiles would characterize where GroupDRO's effect is localized. Two pathways exist; understanding when each operates is the first step toward choosing or combining them wisely.
